% AxiomForge manuscript — template version 1.0.0.
% Build: latexmk -pdf main.tex   (see README.md)
\documentclass[11pt]{amsart}
\makeatletter
% AxiomForge manuscript preamble — template version 1.0.0.
% This file is copied into every contribution and is owned by it afterwards.
% It fixes the typographic identity of the series; it defines no mathematical notation
% (each paper keeps its own notation in macros.tex).

% ---------------------------------------------------------------- page and fonts
\usepackage[a4paper,hmargin=31mm,top=34mm,bottom=36mm,headsep=10mm,footskip=12mm]{geometry}
\usepackage[T1]{fontenc}
\usepackage[utf8]{inputenc}
\usepackage{lmodern}
\usepackage{amsmath,amssymb,amsthm,mathtools}
\usepackage[final,protrusion=true,expansion=true]{microtype}
\usepackage{booktabs}
\usepackage{array}
\usepackage[shortlabels]{enumitem}
\setlist{itemsep=0.2em,topsep=0.4em}
\usepackage{xcolor}
\linespread{1.04}
\setlength{\emergencystretch}{1.5em}
\allowdisplaybreaks[1]
\numberwithin{equation}{section}
\setcounter{tocdepth}{2}
% Table of contents: sections (and their page numbers) in bold with a little space above
% them, subsections indented in normal weight; after \appendix only the (bold) section
% titles are listed, so the appendices never crowd the contents.
\def\l@section{\@tocline{1}{2.5pt plus 1pt minus 2pt}{0pt}{}{\bfseries\let\@tocpagenum\af@tocpagenumbf}}
\def\af@tocpagenumbf#1{\hss{\bfseries #1}}
\def\l@subsection{\@tocline{2}{0pt}{1.5em}{2.3em}{}}
\g@addto@macro\appendix{\addtocontents{toc}{\protect\setcounter{tocdepth}{1}}}
% Every appendix starts on a new page; the unnumbered back matter (\afbacksection) keeps the
% original \section and flows normally after the last appendix.
% (Local implementation through a hook, so that \section itself is not redefined.)
\newif\ifaf@appendixpage
\AddToHook{cmd/section/before}{\ifaf@appendixpage\clearpage\fi}
\g@addto@macro\appendix{\global\af@appendixpagetrue}

% ---------------------------------------------------------------- links
\definecolor{afblue}{rgb}{0.08,0.18,0.42}
\usepackage[pdfusetitle,bookmarks=true,bookmarksnumbered=true,bookmarksopen=true,
  colorlinks=true,linkcolor=afblue,citecolor=afblue,urlcolor=afblue,
  pdfpagelabels=true,pdfdisplaydoctitle=true]{hyperref}
\usepackage[capitalise,noabbrev]{cleveref}
% Local addition: unique PDF anchors for theorem-like environments numbered within sections.
\hypersetup{hypertexnames=false}

% ---------------------------------------------------------------- theorem environments
% thmtools keeps cleveref names correct for environments sharing one counter.
\usepackage{thmtools}
\declaretheorem[style=plain,numberwithin=section,name=Theorem]{theorem}
\declaretheorem[style=plain,sibling=theorem,name=Proposition]{proposition}
\declaretheorem[style=plain,sibling=theorem,name=Lemma]{lemma}
\declaretheorem[style=plain,sibling=theorem,name=Corollary]{corollary}
\declaretheorem[style=plain,numbered=no,name=Theorem]{theorem*}
\declaretheorem[style=definition,sibling=theorem,name=Definition]{definition}
\declaretheorem[style=definition,sibling=theorem,name=Example]{example}
\declaretheorem[style=remark,sibling=theorem,name=Remark]{remark}
\crefname{equation}{}{}
\Crefname{equation}{Equation}{Equations}

% ---------------------------------------------------------------- series metadata
% \afidentifier{AF-0001}, \afversion{1.0.0}, \afdate{25 November 2026},
% \afrepository{https://github.com/axiomforge-math/mathematics/tree/main/results/...},
% \afdisclosure{short AI-assisted research disclosure},
% \afstatus{verification status} (local addition; optional)
\newcommand{\afidentifier}[1]{\gdef\af@identifier{#1}}
\newcommand{\afversion}[1]{\gdef\af@version{#1}}
\newcommand{\afdate}[1]{\gdef\af@date{#1}}
\newcommand{\afrepository}[1]{\gdef\af@repository{#1}}
\newcommand{\afdisclosure}[1]{\gdef\af@disclosure{#1}}
\newcommand{\afcorrespondence}[1]{\gdef\af@correspondence{#1}}
\gdef\af@correspondence{}
% Local addition: verification status of the contribution, printed as a 'Status' row.
\newcommand{\afstatus}[1]{\gdef\af@status{#1}}
\gdef\af@status{}

% The publication block printed under the abstract on the title page.
\newcommand{\afpublicationblock}{%
  \begingroup\footnotesize\noindent
  \begin{tabular}{@{}l@{\quad}p{0.72\textwidth}@{}}
    \textit{Identifier} & \af@identifier\quad(version \af@version, \af@date)\\[1pt]
    \ifx\af@status\@empty\else\textit{Status} & \af@status\\[1pt]\fi
    \textit{Repository} & \url{\af@repository}\\[1pt]
    \ifx\af@correspondence\@empty\else\textit{Correspondence} & \af@correspondence\\[1pt]\fi
    \textit{Disclosure} & \af@disclosure
  \end{tabular}\par
  \endgroup}

% Front matter: title page, then contents, then the body on a fresh page.
\newcommand{\affrontmatter}{%
  \maketitle
  \vspace*{\fill}
  \afpublicationblock
  \thispagestyle{plain}%
  \clearpage
  {\hypersetup{linkcolor=black}\fontsize{9}{10.6}\selectfont\tableofcontents}
  \clearpage}

% Unnumbered back-matter sections (amsart lists starred sections in the contents).
\newcommand{\afbacksection}[1]{\global\af@appendixpagefalse\section*{#1}}

\makeatother
% Notation local to this paper. Define only what materially simplifies the exposition;
% every macro defined here must be used (unused macros are reported by the quality checks).
% This specimen needs no notation of its own.


\title[A specimen of the series]{A specimen of the AxiomForge manuscript series: Fourier's proof of the irrationality of Euler's number}
\author{AxiomForge}
\urladdr{https://github.com/axiomforge-math}
\subjclass[2020]{Primary 11J72; Secondary 00A05}
\keywords{irrationality, typographic specimen, manuscript template}
\date{25 November 2026}

\afidentifier{specimen (not a numbered contribution)}
\afversion{1.0.0}
\afdate{25 November 2026}
\afrepository{https://github.com/axiomforge-math/mathematics/tree/main/docs/manuscript-sample}
\afcorrespondence{\href{mailto:axiomforge.math@gmail.com}{axiomforge.math@gmail.com}}
\afdisclosure{This specimen was prepared with the assistance of an AI system (Claude Opus 5.5). A detailed AI-assisted research disclosure appears at the end.}

\begin{document}

\begin{abstract}
This document is a typographic specimen of the AxiomForge manuscript series. It presents Fourier's classical proof that $e$ is irrational in the format used for every contribution: a title page with publication metadata, a table of contents on its own page, the mathematical development, appendices, and the back matter on reproducibility and AI-assisted research.
\end{abstract}

\affrontmatter

\section{Introduction}\label{sec:introduction}

The number $e=\sum_{k\ge0}1/k!$ is irrational. This fact, due to Euler, admits a
two-line proof found by Fourier, which serves here to display the typographic
conventions of the series rather than to say anything new.

\begin{theorem}\label{thm:main}
The number $e$ is irrational.
\end{theorem}

The proof in \cref{sec:proof} follows the pattern of most irrationality proofs: one
constructs integers which are nonzero by an arithmetic argument and smaller than one in
absolute value by an analytic argument. Here the integers are the scaled tails
$n!\,(e-\sum_{k\le n}1/k!)$; the analytic estimate is the elementary bound of
\cref{lem:tail}, proved in \cref{app:series}.

\section{Proof of the theorem}\label{sec:proof}

This section deduces \cref{thm:main} from a bound on the tails of the exponential
series.

\begin{lemma}\label{lem:tail}
For every integer $n\ge1$,
\begin{equation}\label{eq:tail}
  0<\sum_{k>n}\frac{n!}{k!}<\frac1n .
\end{equation}
\end{lemma}

\begin{proof}[Proof of \cref{thm:main}]
Suppose that $e=a/b$ with positive integers $a$ and $b$, and let $n\ge b$. Then
\[
  N=n!\Bigl(e-\sum_{k=0}^{n}\frac1{k!}\Bigr)=\frac{n!\,a}{b}-\sum_{k=0}^{n}\frac{n!}{k!}
\]
is an integer, because $b$ divides $n!$ and $k!$ divides $n!$ for $k\le n$. By
\eqref{eq:tail}, $0<N<1/n\le1$, which is impossible for an integer.
\end{proof}

\begin{remark}
The same scheme, with integers produced by more elaborate constructions, underlies the
proofs of the irrationality of $\zeta(2)$ and $\zeta(3)$; see \cite{beukers1979note}.
\end{remark}


\appendix
\section{The tail of the exponential series}\label{app:series}

\begin{proof}[Proof of \cref{lem:tail}]
Every term is positive. For $k>n$ we have $k!/n!\ge(n+1)^{k-n}$, hence
\[
  \sum_{k>n}\frac{n!}{k!}\le\sum_{j\ge1}\frac{1}{(n+1)^{j}}=\frac1n ,
\]
with strict inequality because $k!/n!>(n+1)^{k-n}$ as soon as $k\ge n+2$.
\end{proof}


% back matter starts on a new page
\clearpage
\afbacksection{Reproducibility and formal verification}

This specimen contains no computation and no formalization. In a numbered contribution
this section names the Lean declarations of the main theorems, the Lean toolchain and
Mathlib revision, the outcome of the axiom audit, and the commands reproducing every
certificate, and it reports only verification that was actually performed.

\afbacksection{Acknowledgements}

The proof is classical and is reproduced for illustration only.

\afbacksection{AI-assisted research disclosure}

This specimen was written with the assistance of Claude Opus 5.5 (Anthropic), which
drafted the text and the template under the direction of AxiomForge. No new
mathematics is claimed. In a numbered contribution this section describes the role of
each AI system in exploration, construction of arguments, computation, error detection,
manuscript preparation and formalization, the human contributions, and the status of
independent review. The involvement of artificial intelligence does not constitute
independent mathematical certification.

\emph{Responsibility.} AI systems are not authors of this work and cannot take responsibility
for it. AxiomForge takes full responsibility for the entire content of this manuscript,
including every part produced with AI assistance, and for any errors it may contain. Reports
of errors are welcome at \href{mailto:axiomforge.math@gmail.com}{axiomforge.math@gmail.com};
confirmed errors are corrected openly in a new version, recorded in the changelog of this
contribution, with credit to whoever reported them.


\clearpage
\bibliographystyle{amsplain}
\bibliography{references}

\end{document}
